\documentclass[11pt,a4paper]{article}
\usepackage[T1]{fontenc}
\usepackage[utf8]{inputenc}
\usepackage{lmodern}
\usepackage{amsmath,amssymb,amsthm,mathtools,mathrsfs}
\usepackage{graphicx,float,xcolor,multicol}
\usepackage[shortlabels]{enumitem}
\usepackage[margin=27mm]{geometry}
\usepackage{microtype,needspace,placeins}
\usepackage{tikz,tikz-cd}
\usetikzlibrary{arrows.meta,calc,positioning,patterns,decorations.pathreplacing}
\usepackage[colorlinks=true,linkcolor=teal,urlcolor=teal,citecolor=teal]{hyperref}
\setlength{\parindent}{0pt}
\setlength{\parskip}{.6em}
\setcounter{secnumdepth}{0}
\allowdisplaybreaks
\raggedbottom
\providecommand{\cross}{\times}
\DeclareUnicodeCharacter{2020}{\ensuremath{\dagger}}
\DeclareMathOperator{\supp}{supp}
\DeclareMathOperator*{\esssup}{ess\,sup}
\providecommand{\chapter}{\section}
\newenvironment{tool}[2]{\subsection*{Tool #1 --- #2}}{}
\newcommand{\ToolRef}[1]{Tool #1}
\newcommand{\FigureTag}[2]{\textit{Figure #1}}
\newcommand{\Result}[1]{\subsection*{#1}}
\newcommand{\Application}[1]{\subsubsection*{Application\if\relax\detokenize{#1}\relax\else\space--- #1\fi}}
\newcommand{\ProofHeading}{\par\textit{Proof.}\space}
\newcommand{\ProofOf}[1]{\par\textit{Proof of #1.}\space}
\newcommand{\RemarkHeading}{\par\textbf{Remark.}\space}
\newcommand{\NoteHeading}{\par\textbf{Note.}\space}
\newcommand{\ExerciseHeading}{\par\textbf{Exercise.}\space}

% Rendering support only; mathematical text is in each independent post.
\usepackage{booktabs,array,longtable}
\definecolor{HandoutBlue}{HTML}{2F6690}
\definecolor{HandoutNavy}{HTML}{17365D}
\newtheorem{theorem}{Theorem}
\newtheorem{lemma}[theorem]{Lemma}
\newtheorem{proposition}[theorem]{Proposition}
\newtheorem{corollary}[theorem]{Corollary}
\newtheorem{conjecture}[theorem]{Conjecture}
\newtheorem{definition}{Definition}
\newtheorem{example}{Example}
\newtheorem{namedtoolcounter}{Tool}
\newenvironment{namedtool}[1]{\begin{namedtoolcounter}[#1]}{\end{namedtoolcounter}}
\newtheorem*{claim}{Claim}
\newtheorem*{remark}{Remark}
\newtheorem*{exercise}{Exercise}
\newtheorem*{question}{Question}
\newtheorem*{observation}{Observation}
\newcommand{\SourceChapter}[2]{\section*{#2}}
\newcommand{\Lecture}[2]{\section*{Lecture #1: #2}}
\newcommand{\unclear}[1]{\textit{[unclear: #1]}}
\newcommand{\sourcegap}{\textit{[unfinished in the source]}}
\DeclareMathOperator{\dist}{dist}
\DeclareMathOperator{\id}{id}
\DeclareMathOperator{\Hol}{Hol}
\DeclareMathOperator{\Per}{Per}
\DeclareMathOperator{\Fix}{Fix}
\DeclareMathOperator{\diam}{diam}
\DeclareMathOperator{\Var}{Var}
\DeclareMathOperator{\Leb}{Leb}
\DeclareMathOperator{\Hom}{Hom}
\DeclareMathOperator{\End}{End}
\DeclareMathOperator{\Ann}{Ann}
\DeclareMathOperator{\Spec}{Spec}
\DeclareMathOperator{\Max}{Max}
\DeclareMathOperator{\Ob}{Ob}
\DeclareMathOperator{\Tor}{Tor}
\DeclareMathOperator{\Ass}{Ass}
\DeclareMathOperator{\Mor}{Mor}
\DeclareMathOperator{\Fun}{Fun}
\DeclareMathOperator{\Nat}{Nat}
\DeclareMathOperator{\Cone}{Cone}
\DeclareMathOperator{\MCG}{MCG}
\DeclareMathOperator{\Aut}{Aut}
\DeclareMathOperator{\Deck}{Deck}
\DeclareMathOperator{\SL}{SL}
\DeclareMathOperator{\PSL}{PSL}
\DeclareMathOperator{\Area}{Area}
\DeclareMathOperator{\im}{im}
\DeclareMathOperator{\PSh}{PSh}
\newcommand{\CRing}{\mathsf{CRings}}
\newcommand{\AMod}{A\text{-}\mathsf{Mod}}
\newcommand{\Ab}{\mathsf{Ab}}
\newcommand{\Z}{\mathbb Z}
\newcommand{\Q}{\mathbb Q}
\newcommand{\R}{\mathbb R}
\newcommand{\C}{\mathbb C}
\newcommand{\N}{\mathbb N}
\newcommand{\kk}{\mathbb k}
\renewcommand{\aa}{\mathfrak a}
\newcommand{\bb}{\mathfrak b}
\newcommand{\pp}{\mathfrak p}
\newcommand{\qq}{\mathfrak q}
\newcommand{\mm}{\mathfrak m}
\newcommand{\nn}{\mathfrak n}
\newcommand{\rad}{\mathrm r}
\newcommand{\cat}[1]{\mathsf{#1}}
\setcounter{secnumdepth}{0}
\allowdisplaybreaks

\graphicspath{{figures/commutative-algebra/}}
\begin{document}
\section*{Noetherian and Artinian Rings}
% BLOG-CONTENT-BEGIN
\subsection{Noetherian modules}
\setcounter{definition}{10}\begin{definition}
A module $M$ is \emph{Noetherian} (respectively, \emph{Artinian}) if every ascending (respectively, descending) chain of submodules is eventually constant. A ring $A$, viewed as a module over itself, is Noetherian (Artinian) if every ascending (descending) chain of ideals is eventually constant.
\end{definition}
\setcounter{definition}{11}\begin{definition}
For a ring $R$, let
\[d:=\sup\{n\mid\pp_0\subsetneq\pp_1\subsetneq\cdots\subsetneq\pp_n,\ \pp_i\text{ prime in }R\}.\]
Then $d$ is the \emph{dimension} of $R$, written $\dim(R)=d$.
\end{definition}
\paragraph{Comment.}
Every Artinian ring is Noetherian of dimension zero, as proved below; we begin with Noetherian rings.
\setcounter{example}{7}\begin{example}
\begin{enumerate}[label=(\roman*)]
\setcounter{enumi}{0}\item $\Z/n\Z$, as a $\Z$-module, is Artinian and Noetherian.
\setcounter{enumi}{1}\item $\kk[x]$, being a PID, is Noetherian, but $(x)\supset(x^2)\supset(x^3)\supset\cdots$ forces it not to be Artinian.
\setcounter{enumi}{2}\item $\kk[x_1,x_2,\ldots]$ is neither Noetherian nor Artinian.
\setcounter{enumi}{3}\item If $R\hookrightarrow K$, where $K$ is the field of fractions of $R$, then $K$ is both Artinian and Noetherian, but $R$ can be neither, as in (iii), which is merely a domain.
\end{enumerate}
\end{example}
\setcounter{theorem}{12}\begin{proposition}
$M$ is Noetherian if and only if every submodule of $M$ is finitely generated.
\end{proposition}
\begin{proof}
Suppose $M$ is Noetherian and $N$ is a submodule of $M$. Set
\[\Sigma=\{N_0\subseteq N\mid N_0\text{ is a finitely generated submodule}\}.\]
Since $0\in\Sigma$, this set is nonempty. For every chain in $\Sigma$, its union is a submodule of $M$ and is finitely generated.
% Source page 2
Its finitely many generators all lie in one member of the chain, so the union equals that member and belongs to $\Sigma$. Zorn's lemma gives a maximal $N_0\in\Sigma$. If $N_0\ne N$, choose $x\in N\setminus N_0$; then $N_0\subset N_0+(x)$, absurd.

Conversely, let $M_1\subseteq M_2\subseteq\cdots$ be a chain of submodules. Then $\bigcup_{i=1}^{\infty}M_i$ is a submodule of $M$, hence is finitely generated, say by $x_1,\ldots,x_n$. If $x_i\in M_{n_i}$, then for $n=\max_i\{n_i\}$ one has $M_n=\bigcup_{i=1}^{\infty}M_i$.
\end{proof}
\setcounter{theorem}{13}\begin{proposition}
If $0\to M'\to M\xrightarrow{\beta}M''\to0$ is exact, then $M$ is Noetherian if and only if $M'$ and $M''$ are Noetherian.
\end{proposition}
\begin{proof}
If $M$ is Noetherian, every chain in $M'$ is a chain in $M$, and a chain in $M''$ corresponds to a chain in $M$ containing $M'$.

Conversely, let $L$ be a submodule of $M$. Let $\{x_i\}_{i=1}^n$ generate $\beta(L)$ and $\{y_j\}_{j=1}^m$ generate $L\cap M'$. For $x\in L$, write
\[\beta(x)=\sum_{i=1}^n r_i x_i\qquad(r_i\in A,\ x_i\in M'').\]
Since $M/M'\cong M''$, write $x_i=[\widetilde x_i]$, with $\widetilde x_i\in L$, and $\beta(x)=[x]$. Then
\[
x\equiv\sum_{i=1}^n r_i\widetilde x_i\pmod{M'},\qquad
x-\sum_{i=1}^n r_i\widetilde x_i\in L\cap M'.
\]
Thus $x-\sum_{i=1}^n r_i\widetilde x_i=\sum_{j=1}^m s_jy_j$, so
\[x=\sum_{i=1}^n r_i\widetilde x_i+\sum_{j=1}^m s_jy_j\qquad(x\in L).\]
Therefore $\{\widetilde x_i\}_{i=1}^n\cup\{y_j\}_{j=1}^m$ generates $L$, and $M$ is Noetherian.
\end{proof}
\setcounter{theorem}{14}\begin{corollary}
\begin{enumerate}[label=(\roman*)]
\setcounter{enumi}{0}\item If $M_1,\ldots,M_n$ are Noetherian, then $\bigoplus_{i=1}^n M_i$ is Noetherian.
\setcounter{enumi}{1}\item If $M$ is finitely generated over Noetherian $A$, then $M\cong A^n/P$ is Noetherian.
\setcounter{enumi}{2}\item If $\aa<A$ is an ideal of a Noetherian ring, then $A/\aa$ is Noetherian.
% Source page 3
\setcounter{enumi}{3}\item If $\varphi:A\twoheadrightarrow B$ and $A$ is Noetherian, then $B\cong A/\ker\varphi$ is Noetherian.
\setcounter{enumi}{4}\item If $A$ is Noetherian, then $S^{-1}A$ is Noetherian. In particular, for prime $\pp$, $A_{\pp}$ is Noetherian.
\end{enumerate}
\end{corollary}
\setcounter{theorem}{15}\begin{theorem}[Hilbert's basis theorem]
If $A$ is Noetherian, then $A[x]$ is Noetherian.
\end{theorem}
\begin{proof}
Let $\aa<A[x]$ be an ideal. Construct $I=(\text{leading coefficients of elements in }\aa)$. Since $I<A$, it is finitely generated, say by $\{a_i\}_{i=1}^n$. For every $1\le i\le n$, there exists $f_i\in\aa$ with $f_i=a_ix^{r_i}+\text{lower terms}$. Set $r=\max_{1\le i\le n}\{r_i\}$, and let $\aa':=(f_1,\ldots,f_n)\subset\aa$.

Let $f=ax^m+\text{lower terms}\in\aa$. Then $a\in I$, so $a=\sum_{i=1}^n u_ia_i$. If $m\ge r$, then $f-\sum_{i=1}^n u_if_ix^{m-r_i}\in\aa$ has degree less than $m$. Proceeding in this way gives $f=g+h$, with $\deg g<r$ and $h\in\aa'$.

Let $M$ be the $A$-submodule of $A[x]$ generated by $1,x,\ldots,x^{r-1}$. Then $\aa=(\aa\cap M)+\aa'$. Since $A$ is Noetherian, the finitely generated module $M$ is Noetherian, hence $\aa\cap M$ is finitely generated. If $\{g_i\}_{i=1}^k$ generates $\aa\cap M$, then $\{f_i\}_{i=1}^n\cup\{g_i\}_{i=1}^k$ generates $\aa$.
\end{proof}
\setcounter{theorem}{16}\begin{corollary}
\begin{enumerate}[label=(\roman*)]
\setcounter{enumi}{0}\item If $A$ is Noetherian, then $A[x_1,\ldots,x_n]$ is Noetherian.
% Source page 4
\setcounter{enumi}{1}\item A finitely generated algebra over Noetherian $A$ is Noetherian.
\end{enumerate}
\end{corollary}
\subsection{Noether's work on primary decomposition}
\setcounter{definition}{12}\begin{definition}
An ideal $\aa<A$ is \emph{irreducible} if $\aa=\bb\cap\mathfrak c\Rightarrow\aa=\bb$ or $\aa=\mathfrak c$.
\end{definition}
\setcounter{theorem}{17}\begin{proposition}
Every irreducible ideal in a Noetherian ring is primary.
\end{proposition}
\begin{proof}
Let $\aa<A$ be irreducible. Then $\aa$ is primary if and only if $(0)$ is primary in the Noetherian ring $A/\aa$. Let $xy=0$, $y\ne0$. Consider $\Ann(x)\subseteq\Ann(x^2)\subseteq\cdots$. There exists $n$ such that $\Ann(x^n)=\Ann(x^{n+1})=\cdots$.

\emph{Claim:} $(x^n)\cap(y)=(0)$. If $a\in(y)$, then $ax=0$; and if $a\in(x^n)$, then $a=bx^n$, so $bx^{n+1}=0$. Hence $b\in\Ann(x^{n+1})=\Ann(x^n)$, giving $a=0$.

Since $\aa$ is irreducible, so is $(0)$ in $A/\aa$. As $y\ne0$, we get $x^n=0$. Thus $(0)$ is primary in $A/\aa$.
\end{proof}
\setcounter{theorem}{18}\begin{theorem}
Every ideal in a Noetherian ring admits primary decomposition.
\end{theorem}
\begin{proof}
Set $\Sigma=\{\aa<A\mid\aa\text{ is not a finite intersection of irreducible ideals}\}$. If $\Sigma\ne\varnothing$, then, as every chain in $A$ stabilizes, $\Sigma$ satisfies the hypotheses of Zorn's
% Source page 5
lemma, which yields a maximal element $\aa$. This ideal is reducible, so $\aa=\bb\cap\mathfrak c$ with $\aa\subset\bb$ and $\aa\subset\mathfrak c$. By maximality, $\bb,\mathfrak c\notin\Sigma$; they are finite intersections of irreducible ideals, and hence so is $\aa$. This is absurd. Therefore $\Sigma=\varnothing$.
\end{proof}
\subsection{Artinian rings}
The following propositions show that Artinian rings are Noetherian of dimension zero.
\begin{enumerate}[label=(\roman*)]
\setcounter{enumi}{0}\item In an Artinian ring every prime is maximal. If $\pp$ is prime, then $A/\pp$ is an Artinian integral domain. For $0\ne x\in A/\pp$, there is some $n$ with $(x^n)=(x^{n+1})$. Thus $x^n=x^{n+1}y$. Since $A/\pp$ is a domain, $x^n\ne0$ and $xy=1$. Hence $A/\pp$ is a field and $\pp$ is maximal.
\setcounter{enumi}{1}\item An Artinian ring has only finitely many maximal ideals. Define
\[\Sigma=\left\{\bigcap_{i=1}^n\mm_i\ \middle|\ n\in\N,\ \mm_i\text{ maximal}\right\}.\]
Since $A$ is Artinian, every chain in $\Sigma$ has a lower bound. By Zorn's lemma, there exists a minimal $\bigcap_{i=1}^{n_0}\mm_i\in\Sigma$. For any maximal ideal $\mm$, minimality yields $\mm\cap(\bigcap_i\mm_i)=\bigcap_i\mm_i$.
% Source page 6
Thus $\bigcap_i\mm_i\subseteq\mm$, so $\mm_i=\mm$ for some $i$.
\setcounter{enumi}{2}\item The nilradical of an Artinian ring is nilpotent.
\end{enumerate}
\setcounter{theorem}{19}\begin{theorem}
If $A$ is Artinian, then $A$ is Noetherian and $\dim(A)=0$.
\end{theorem}
\begin{proof}[Sketch]
If $A$ is Artinian, then $\dim(A)=0$. Let $\{\mm_i\}_{i=1}^n$ be the maximal ideals of $A$. For some $k$,
\[
\prod_{i=1}^n\mm_i^k\subseteq\left(\bigcap_{i=1}^n\mm_i\right)^k
=\bigl(\rad((0))\bigr)^k=0.
\]
\emph{Fact:} If $\prod_{i=1}^n\mm_i=0$, with the $\mm_i$ maximal and not necessarily distinct, then $A$ is Noetherian.

Indeed, $A$ is Noetherian if and only if $\mm_1$ and $A/\mm_1$ are Noetherian. The latter is a field, hence Noetherian, so $A$ is Noetherian if and only if $\mm_1$ is Noetherian. Further, $\mm_1$ is Noetherian if and only if $\mm_1\mm_2$ and $\mm_1/\mm_1\mm_2$ are Noetherian. Since $\mm_2(\mm_1/\mm_1\mm_2)=0$, the quotient is an $A/\mm_2$-vector space. The module $\mm_1$ is Artinian because $A$ is Artinian; hence $\mm_1/\mm_1\mm_2$ is an Artinian $A/\mm_2$-module, thus a finite-dimensional vector space. It is therefore a Noetherian $A/\mm_2$-module, and consequently a Noetherian $A$-module. Thus $\mm_1$ is Noetherian if and only if $\mm_1\mm_2$ is Noetherian.

Proceeding inductively to $\mm_1\mm_2\cdots\mm_n=0$ shows that $A$ is Noetherian.
\end{proof}
% Source page 7
\begin{remark}
The converse is also true.
\end{remark}
% BLOG-CONTENT-END
\end{document}
